\begin{textsample}{NEG Dim 4 – 2022\_c5 – Score -1.00 – 2022\_c5/enade\_2022\_tec\_gest\_tur\_eadd874aa1.txt}  \label{ex:f4_neg_080}
CAPÍTULO 6 DISTRIBUIÇÃO DE CONCEITOS A finalidade deste capítulo é apresentar um panorama nacional sobre a distribuição do Conceito Enade 2022.
Para traçar esse panorama recorre -se a tabelas e gráficos, que são acompanhados de comentários e análises que contribuem para a contextualização dos dados neles expostos.
Para visualizar o panorama do Enade/2022, em nível nacional, na primeira seção do capítulo, os dados sobre o número de cursos participantes são organizados por Grande região.
Na segunda, o mapeamento do conceito do exame é feito por categoria administrativa das IES e modalidade de oferta dos cursos, de maneira cruzada às Grandes regiões.
Na última seção, a lógica de apresentação dos dados é feita por organização acadêmica e, também, Grande região de funcionamento dos cursos, no país como um todo. 6.1 PANORAMA NACIONAL DA DISTRIBUIÇÃO DOS CONCEITOS A lógica adotada na apresentação dos dados, nesta seção e nas duas seguintes, pressupõe o registro de duas informações de natureza reiterativa.
A primeira informação diz respeito aos valores possíveis do Enade, que variam em níveis de 1 a 5, sendo que, à medida que esses valores aumentam, o curso recebe uma avaliação melhor ; além disso, os cursos com somente um ou sem nenhum concluinte participante não obtêm o Conceito Enade, ficando Sem Conceito ( SC ), como já ressaltado neste relatório.
O segundo registro é que a diferença entre os cursos tabulados neste capítulo e no Capítulo 3 corresponde aos cursos que não foram avaliados, ou seja, àqueles, em princípio, sem estudantes concluintes que tivessem sido inscritos para a prova.
Na Tabela 6.1, são apresentadas a quantidade e a distribuição de cursos de Turismo participantes do Enade/2022, por faixa de conceito e Grande região.
Observando -se os dados da Tabela 6.1, nota -se que, dos 83 cursos participantes, 22 ( 26,5 \% ) classificaram -se com conceito 3, o valor modal.
Esse foi também o conceito modal na maioria das regiões : Norte ( 40 \% ), Sul ( 26,7 \% ) e Centro-Oeste ( 44,4 \% ).
O conceito 2 foi o segundo mais frequente em nível nacional ( 25,3 \%, correspondendo a 21 cursos ), tendo sido o conceito modal na região Sudeste ( 25,7 \% ) e na Nordeste ( 26,3 \% ).
O conceito 4 foi o terceiro ( 22,9 \%, correspondendo a 19 cursos ).
Houve, ainda, 9 cursos ( 10,8 \% ) que receberam o conceito 5, e 4 cursos ( 4,8 \% ) que receberam o conceito 1.
Dos 83 cursos de Turismo, 8 ( 9,6 \% ) ficaram Sem Conceito ( SC ).
A região Norte participou com cinco cursos, o que corresponde a 5,5 \% do total nacional.
Desses, dois cursos receberam o conceito 3, o que equivale a 40 \% do total regional, o conceito modal na região, como já comentado.
A um curso ( 20 \% ) atribuiu -se o conceito 1, e a um curso ( 20 \% ), o conceito 2.
Nenhum curso ficou Sem Conceito ( SC ), assim como o conceito 5, e um curso ( 20 \% ), com o conceito 4, nessa região.
A região Nordeste participou com 19 cursos, o correspondente a 19,1 \% do total nacional.
Nessa região, seis cursos ( 31,6 \% em termos regionais ) obtiveram conceito 4, o conceito modal para a região.
Os conceitos 1 e 2 foram atribuídos, respectivamente, a um e a cinco cursos ( 5,3 \% e 26,3 \%, respectivamente ).
Ainda nessa região, quatro cursos ( 21,1 \% ) foram avaliados com o conceito 3, e dois cursos ( 10,5 \% ) foram avaliados com o conceito 5.
Um curso ficou Sem Conceito ( SC ).
Para os 35 cursos participantes da região Sudeste ( 42 \% do total nacional ), o conceito 2 foi o conceito modal, como já anteriormente assinalado, sendo atribuído a nove cursos ( 25,7 \% ).
O conceito 1 foi atribuído a um curso ( 2,9 \% ), e o conceito 3, a oito cursos ( 22,9 \% ).
Receberam os conceitos 4 e 5, respectivamente, nove cursos ( 25,7 \% ) e quatro cursos ( 11,4 \% ).
Nessa região, quatro cursos ficaram Sem Conceito ( SC ).
Os 15 cursos da região Sul corresponderam a 18 \% do total nacional.
Como já indicado, houve predominância do conceito 3 ( conceito modal ), atribuído a quatro ( 26,7 \% ) cursos dessa região.
Ao conceito 1 não foram atribuídos cursos, ao conceito 2, foram atribuídos três cursos ( 20 \% ) ; e aos conceitos 4 e 5, atribuídos a três ( 20 \% ) e a dois ( 13,3 \% ) cursos, respectivamente.
Na região Sul, três cursos ficaram Sem Conceito ( SC ).
Já dos nove cursos participantes na região Centro-Oeste ( 11,0 \% do total nacional ), quatro ( 44,4 \% em termos regionais ) receberam o conceito 3, o conceito modal, como já destacado.
Um curso ( 11,1 \% ) obteve o conceito 1, 3 cursos ( 33,3 \% ) obtiveram o conceito 2 ; nenhum curso obteve o conceito 4, e um curso ( 11,1 \% ), o conceito 5.
Nessa região, nenhum curso ficou Sem Conceito ( SC ).
Em um gráfico de distribuição cumulativa, em cada valor de abscissa ( eixo x ), o valor apresentado no eixo das ordenadas ( eixo y ) é o valor acumulado da variável ( Conceito Enade ).
Em linhas gerais, em um gráfico desta mesma tipologia, poligonais mais à esquerda representam, no caso dos dados apresentados nesta seção, as grandes regiões com a pior distribuição, e poligonais mais à direita, as grandes regiões com os melhores conceitos.
A partir destas diretrizes, no Gráfico 6.1, é apresentada a distribuição cumulativa do Conceito Enade, segundo a grande região.
Nesse gráfico, poligonais mais à esquerda representam valores mais baixos do Conceito Enade, e poligonais mais à direita, valores mais altos. 6.2 CONCEITOS POR CATEGORIA ADMINISTRATIVA, MODALIDADE DE OFERTA DOS CURSOS E GRANDE REGIÃO Na Tabela 6.2, é apresentada a distribuição dos cursos de Turismo participantes do Enade/2022, por categoria administrativa e modalidade de oferta, de acordo com os conceitos por eles alcançados, segundo as grandes regiões brasileiras.
De acordo com as informações da Tabela 6.2, dos 83 cursos participantes, 49 ( 59 \% ) eram ministrados em Instituições Públicas e 34 ( 41 \% ), em Instituições Privadas.
Quanto à modalidade de oferta, quatro cursos eram a Distância e 79, presenciais.
Ainda conforme as informações da mesma tabela, em termos nacionais, entre instituições públicas e privadas, observa -se que, dos nove cursos avaliados com conceito 5, oito foram oferecidos em IES públicas, e 1 um em IES privadas.
Dos 49 cursos participantes de IES públicas, o conceito 4 foi o valor modal, atribuído a 16 cursos ( 32,6 \% ).
Dos demais cursos de instituições públicas participantes, um curso ( 2 \% ) foi avaliado com o conceito 1, 11 cursos ( 22,4 \% ) foram avaliados com o conceito 2, 13 cursos ( 26,5 \% ) foram avaliados com o conceito 3, e oito cursos ( 16,3 \% ), com o conceito 5.
Nessa categoria, nenhum curso ( 0 \% ) ficou Sem Conceito ( SC ).
Na rede privada, o conceito modal foi o 2, atribuído a dez cursos, o correspondente a 29,4 \% dos 34 cursos da categoria.
Dos demais cursos participantes, três ( 8,8 \% ) receberam o conceito 1, e nove ( 26,5 \% ), o conceito 3.
O conceito 4 foi atribuído a três cursos ( 8,8 \% ), e o conceito 5, a um curso ( 2,9 \% ).
Nessa categoria administrativa, oito cursos ( 23,6 \% ) ficaram Sem Conceito ( SC ).
Dos 79 cursos participantes da modalidade presencial, o conceito 3 foi o valor modal, atribuído a 20 cursos ( 25,3 \% ).
Dos demais cursos nessa modalidade de oferta, quatro cursos ( 5 \% ) receberam o conceito 1 ; 19 cursos ( 24 \% ) receberam o conceito 2 ; 19 cursos ( 24 \% ), o conceito 4 ; e nove cursos ( 11,4 \% ), o conceito 5.
Nessa modalidade de oferta, oito cursos ( 10,1 \% ) ficaram Sem Conceito ( SC ).
Dos quatro cursos participantes da modalidade a distância, os conceitos 2 e 3 foram os valores bimodais, atribuídos a dois cursos ( 50 \% ).
Nessa modalidade de oferta, nenhum curso ficou Sem Conceito ( SC ).
Na análise por região, observa -se que, na região Norte, as instituições privadas participaram com dois cursos ( 40 \% do total regional ), e nenhum curso ficou Sem Conceito ( SC ).
Os conceitos modais para as instituições privadas na região foram o 1 e o 2.
As instituições públicas participaram com três cursos na região Norte ( 60 \% do total regional ), e nenhum curso ficou Sem Conceito ( SC ).
O conceito modal nas instituições públicas na região foi o 3, atribuído a dois cursos, o correspondente a 66,6 \% dessa categoria na região.
Um curso foi avaliado com conceito 4.
Dos cursos oferecidos na região Norte, cinco eram presenciais, e nenhum ficou Sem Conceito ( SC ).
Dos cinco cursos oferecidos, dois obtiveram o conceito 3 ; e cada um dos demais, os conceitos 1, 2 e 4.
Nenhum curso ficou Sem Conceito ( SC ).
Na região Nordeste, a rede privada concentrou três dos 19 cursos participantes, o equivalente a 15,8 \% do total da região.
Entre os cursos oferecidos por IES privadas no Nordeste, um curso foi avaliado com os conceitos 2 e 3.
Nessa Categoria administrativa e Grande região, um curso ficou Sem Conceito ( SC ).
As instituições públicas da região Nordeste participaram com 16 cursos ( 84,2 \% do total da região ).
Desses, 6 ( 37,5 \% ) obtiveram conceito 4, o valor modal, seguidos por quatro cursos, que receberam o conceito 2 ; três, o conceito 3 ; dois, o conceito 5 e um, o conceito 1.
Nenhum curso ficou Sem Conceito ( SC ).
Dos cursos oferecidos na região Nordeste, 17 eram presenciais : um ficou Sem conceito ( SC ), um ficou com o conceito 1 ; 4, com o conceito 2 ; 3, com o conceito 3 ; 6, com o conceito 4 ( sendo o valor modal ) e dois, com o conceito 5.
Na modalidade a distância, foram 2 cursos, tendo um obtido os conceitos 2 e 3.
Nessa região, nenhum curso ficou Sem Conceito ( SC ) ou recebeu os conceitos 1, 4 e 5.
Na região Sudeste, o percentual de cursos da rede privada, 60 \%, correspondeu a 21 dos 35 cursos participantes.
Nessa combinação de Categoria administrativa e grande região, o conceito bimodal foi o 2 e o 3, atribuído a sete cursos.
Quatro cursos ficaram Sem Conceito ( SC ).
Os demais foram avaliados com os seguintes conceitos : 1 ( um curso ) e 4 ( dois cursos ).
Dos 14 cursos oferecidos por instituições públicas na região Sudeste, o conceito modal foi o 4, atribuído a sete cursos.
Dos demais cursos, dois receberam o conceito 2 ; um recebeu o conceito 3 ; e quatro, o conceito 5.
Nessa combinação de categoria administrativa e grande região, nenhum curso ficou Sem Conceito ( SC ), e nenhum recebeu o conceito 1.
Também na região Sudeste, grande parcela dos cursos foi oferecida na modalidade presencial : 34 de 35 cursos.
Para esses, a bimodal foram os conceitos 2 e 3, com 8 cursos cada.
Os demais cursos presenciais da região foram \textbf{alocados} nos conceitos 1 ( um curso ), 4 ( 9 cursos ) e 5 ( 4 cursos ).
Sendo que quatro cursos presenciais ficaram Sem Conceito ( SC ) nessa região.
Ainda na região Sudeste, somente 1 curso foi oferecido na modalidade a distância, tendo conceito 2.
E nenhum curso a distância ficou Sem Conceito ( SC ).
As instituições privadas concentraram sete dos 15 cursos participantes na região Sul, 46,7 \% do total regional.
Desses, um curso obteve conceito 2, 3, 4 e 5.
Nessa combinação de categoria administrativa e grande região, três cursos ficaram Sem Conceito ( SC ).
As instituições públicas na região Sul participaram com oito cursos ( 53,3 \% dos cursos da região ), aos quais foram atribuídos os conceitos 3 ( três cursos ), 2 e 4 ( dois cursos cada ) e 5 ( um curso ).
Nenhum curso da região Sul oferecido por IES públicas ficou Sem Conceito ( SC ) ou recebeu o conceito 1.
Na região Sul, 14 dos 15 cursos foram oferecidos na modalidade presencial.
Para esses, a trimodal foram os conceitos 2, 3 e 4, com três cursos cada.
Os demais cursos presenciais da região foram \textbf{alocados} no conceito 5 ( dois cursos ) e três cursos ficaram Sem Conceito ( SC ).
Na modalidade a distância, foi oferecido somente um curso, com conceito 3.
Nenhum curso ficou Sem Conceito ( SC ) ou com o conceito 1, 2, 4 e 5.
Na região Centro-Oeste, um dos nove cursos participantes eram de IES privadas ( 11,1 \% em termos regionais ), recebeu conceito 1.
Nenhum curso ficou Sem Conceito ( SC ).
Quanto aos oito cursos oferecidos por instituições públicas na região Centro-Oeste ( 88,9 \% do total regional ), quatro cursos apresentaram o conceito 3, conceito modal ; três cursos obtiveram o conceito 2 e 1, o conceito 5.
Nenhum curso ficou Sem Conceito ( SC ) ou com o conceito 4.
Ainda na região Centro-Oeste, os nove ( 100 \% ) cursos foram presenciais.
Para esses, a moda foi o conceito 3, com quatro cursos ( 44,4 \% ).
Os demais cursos presenciais da região foram \textbf{alocados} nos conceitos 1 ( um curso, 11,1 \% ), 2 ( três cursos, 33,3 \% ) e 5 ( 1 curso, 11,1 \% ).
Nenhum curso presencial da região ficou Sem Conceito ( SC ) ou com o conceito 4.
À título de síntese, no Gráfico 6.2, apresenta -se a distribuição do Conceito Enade, segundo a categoria administrativa da IES.
Nesse gráfico, como dito anteriormente, poligonais mais à esquerda representam valores mais baixos do Conceito Enade, e poligonais mais à direita, valores mais altos.
É importante observar que, para cada poligonal presente nesse gráfico, a porcentagem constante no eixo y se refere ao universo de cursos oferecidos por instituições públicas ( poligonal vermelha ) ou privadas ( poligonal azul ). 6.3 CONCEITOS POR ORGANIZAÇÃO ACADÊMICA E GRANDE REGIÃO Na Tabela 6.3, encontra -se a distribuição dos conceitos atribuídos aos cursos participantes do Enade/2022, na área de Turismo, por organização acadêmica, segundo as grandes regiões brasileiras.
Como mostram os dados da Tabela 6.3, dos 83 cursos de Turismo participantes, 64 eram oferecidos em universidades, seis em Centros Universitários, oito em Faculdades e cinco em CEFET/IF.
Esta distribuição corresponde a, respectivamente, 77,1 \%, 7,2 \%, 9,6 \% e 6,02 \% dos cursos.
De acordo com os dados apresentados na mesma tabela, dos nove cursos avaliados com o conceito 5 estavam vinculados a universidades.
Os cursos vinculados a esse tipo de organização acadêmica tiveram conceitos 2 e 3, com o conceito bimodal, atribuídos a 17 cursos ( 20,5 \% ) cada um.
Os demais cursos ofertados por universidades receberam os conceitos 1 ( três cursos ), 4 ( 16 cursos ) e 5 ( nove cursos, como já mencionado ).
Dois cursos ficaram Sem Conceito ( SC ).
Dos cursos vinculados a Centros Universitários, o conceito modal também foi o 3, atribuídos a três cursos ( 50 \% ).
Dois cursos vinculados a esse tipo de Organização acadêmica ficaram Sem Conceito ( SC ), e os outros cursos receberam o conceito 2 ( 1 curso ).
Dos oito cursos mantidos por Faculdades, dois ( 25 \% ) receberam conceito 2, o conceito modal.
O segundo conceito mais frequentemente atribuído aos cursos oferecidos por essa categoria de IES foram os conceitos 1 e 3, recebido por um curso cada.
Dos cursos oferecidos em Faculdades, quatro ficaram Sem Conceito ( SC ).
Dos 5 cursos mantidos por CEFET/IF, 3 ( 60 \% ) receberam conceito 4, o conceito modal.
O segundo conceito mais frequentemente atribuído aos cursos oferecidos por essa categoria de IES foram os conceitos 2 e 3, recebidos por um curso cada.
Nenhum curso ficou Sem Conceito ( SC ) ou com os conceitos 1 e 5.
Considerando -se, separadamente, as regiões brasileiras, verifica -se que, na região Norte, as universidades participaram com quatro cursos avaliados, aos quais foram atribuídos os seguintes conceitos : 1 ( a um curso ) ; 4 ( a um curso ) ; 3, o conceito modal ( a dois cursos ).
Nenhum curso ficou Sem Conceito ( SC ) nessa combinação de categoria e região.
Os Centros Universitários da região Norte foram representados por um curso apenas, com conceito 2.
Nenhum curso ficou Sem Conceito ( SC ) ou com os conceitos 1, 3, 4 e 5.
Na região Nordeste, as universidades participaram com 15 dos 19 cursos da Área de Turismo oferecidos.
O conceito modal foi 4, atribuído a cinco cursos.
Os quatro cursos obtiveram o conceito 2 ; três cursos, o conceito 3 ; e outros dois, o conceito 5.
Nenhum curso oferecido por universidades no Nordeste ficou Sem Conceito ( SC ).
Os Centros Universitários contaram com apenas um curso participante na região Nordeste, sendo um Sem Conceito ( SC ).
As Faculdades foram representadas por dois cursos nessa região, cada um recebeu conceito 2 e 3.
Nenhum curso ficou Sem Conceito ( SC ).
Nessa combinação de região e Organização acadêmica, apenas um curso foi oferecido em CEFET/IF, com conceito 4.
Na região Sudeste, as universidades concentraram 26 dos 35 cursos de Turismo da região.
Dos cursos oferecidos em IES com este tipo de organização acadêmica, o conceito modal foi 2, atribuído a oito cursos.
Os demais cursos receberam os conceitos 4 ( sete cursos ), 3 ( cinco cursos ), 5 ( quatro cursos ) e 1 ( um curso ).
Um curso ficou Sem Conceito ( SC ).
Os Centros Universitários participaram com quatro cursos na região Sudeste, dos quais três obtiveram o conceito modal 3.
Um curso ficou Sem Conceito ( SC ) nessa combinação de organização acadêmica e região.
Já as Faculdades foram representadas por três cursos na região Sudeste, dos quais dois ficaram Sem Conceito ( SC ).
O conceito modal 2 foi atribuído a um curso.
Para os dois cursos participantes de CEFET/IF, na região Sudeste, o conceito foi 4 ( dois cursos ).
Nenhum curso ficou Sem Conceito ( SC ) ou com os conceitos 1, 2, 3 e 5 nesse tipo de organização acadêmica.
Dos 15 cursos da região Sul, 13 eram mantidos por universidades, para os quais o conceito modal foi 3, atribuído a 4 cursos.
Os demais cursos receberam os conceitos 2 ( 3 cursos ), 4 ( 3 cursos ) e 5 ( dois cursos ) ; nenhum curso recebeu o conceito 1.
Em IES com esse tipo de organização acadêmica, um curso da região Sul ficou Sem Conceito ( SC ).
Dos dois cursos participantes das Faculdades da região Sul, duas ficaram Sem Conceito ( SC ).
Na região Centro-Oeste, seis dos nove cursos oferecidos eram mantidos por universidades.
Para esse tipo de Organização acadêmica, o conceito modal foi 3, atribuído a três cursos ; e nenhum curso da região ficou Sem Conceito ( SC ).
Os demais cursos receberam os conceitos 2 ( dois cursos ) e 5 ( um curso ).
Nos Centros Universitários da região Centro-Oeste, não houve cursos.
O único curso oferecido em Faculdades na região Centro-Oeste, recebe o conceito 1.
Nessa região, nenhum curso oferecido em Faculdades ficou Sem Conceito ( SC ) ou com os conceitos 2, 3, 4 e 5.
Dos dois cursos oferecidos em CEFET/IF, um curso foi avaliado com o conceito 2 e um curso obteve o conceito 3.
Nesse tipo de organização acadêmica, nenhum curso da região Centro-Oeste ficou Sem Conceito ( SC ) ou com os conceitos 1, 4 e 5.
De maneira semelhante à forma como se procedeu, em relação à seção anterior, encerra -se esta seção com o Gráfico 6.3, que apresenta a distribuição do Conceito Enade dos cursos de Turismo segundo a organização acadêmica da IES.
A leitura desse último gráfico é semelhante à leitura feita dos anteriores, seguindo -se a lógica de que poligonais deslocadas para a esquerda correspondem a valores mais baixos do Conceito Enade ( 1, 2 e 3 ) e poligonais deslocadas para a direita correspondem a valores mais altos do Conceito Enade.

% matched lemmas: alocar
\end{textsample}
